\heading{数据结构与算法B-24春-期末}
\noteinfo{题目来源：2024 春季期末扫描版，卷面标注为数据结构与算法B。已按扫描件复核选择题、判断题、填空题和算法填空中的可辨认文字；个别原题信息不足处保留说明。本次整理提供 C 语言版与 Python 语言版；除程序代码语言外，两版题目内容保持一致。}

\textbf{一、选择题（30 分，每小题 2 分）}

\begin{question}
设栈的输入序列为 $1,2,3,4,5$，下列序列中不可能是其出栈序列的是（\quad）。
\begin{enumerate}
  \item $23415$
  \item $54132$
  \item $23145$
  \item $15432$
\end{enumerate}
\begin{solution}
选 B。若 $5$ 最先出栈，则 $1,2,3,4$ 已依次在栈内，之后必须按 $4,3,2,1$ 的相对次序出栈，不可能出现 $5,4,1,3,2$。
\end{solution}
\end{question}

\begin{question}
对线性表进行二分法查找，其前提条件是（\quad）。
\begin{enumerate}
  \item 线性表以链接方式存储，并且按关键码值排序
  \item 线性表以链接方式存储，并且按关键码值的检索频率排序
  \item 线性表以顺序方式存储，并且按关键码值排序
  \item 线性表以顺序方式存储，并且按关键码值的检索频率排序
\end{enumerate}
\begin{solution}
选 C。二分查找需要能随机访问中间元素，同时要求关键码有序。
\end{solution}
\end{question}

\begin{question}
对于二叉树中的结点 $N$，若其左子树高度为 $h_1$，右子树高度为 $h_2$，则以 $N$ 为根结点的子树高度为（\quad）。
\begin{enumerate}
  \item $\max(h_1,h_2)$
  \item $\min(h_1,h_2)$
  \item $\max(h_1,h_2)+1$
  \item $\min(h_1,h_2)+1$
\end{enumerate}
\begin{solution}
选 C。根结点所在层还需计入高度。
\end{solution}
\end{question}

\begin{question}
在插入排序算法中，如果待排序序列已经是有序的，则插入排序算法的时间复杂度为（\quad）。
\begin{enumerate}
  \item $O(n^2)$
  \item $O(n\log n)$
  \item $O(n)$
  \item $O(1)$
\end{enumerate}
\begin{solution}
选 C。此时每一趟只需一次比较即可确定无需移动，总时间为线性阶。
\end{solution}
\end{question}

\begin{question}
下列概念中属于存储结构的是（\quad）。
\begin{enumerate}
  \item 线性表
  \item 链表
  \item 栈
  \item 队列
\end{enumerate}
\begin{solution}
选 B。链表是链式存储结构，其余通常为逻辑结构或抽象数据类型。
\end{solution}
\end{question}

\begin{question}
在二分查找算法中，每次比较后都将搜索范围缩小一半。若要在长度为 $n$ 的数组中查找一个元素，最坏情况下需要比较多少次？（\quad）
\begin{enumerate}
  \item $n$
  \item $\lfloor n/2\rfloor$
  \item $\lfloor\log_2 n\rfloor$
  \item $\lfloor\log_2 n\rfloor+1$
\end{enumerate}
\begin{solution}
选 D。二分查找判定树高度为 $\lfloor\log_2 n\rfloor+1$。
\end{solution}
\end{question}

\begin{question}
在一个具有 $n$ 个结点的有序单链表中插入一个新结点并仍保持其有序，平均时间复杂度为（\quad）。
\begin{enumerate}
  \item $O(n\log n)$
  \item $O(1)$
  \item $O(n)$
  \item $O(n^2)$
\end{enumerate}
\begin{solution}
选 C。需要顺链查找插入位置，平均需扫描线性数量的结点。
\end{solution}
\end{question}

\begin{question}
设有 $4$ 棵结点关键码均为整数的二叉树，若其中序遍历序列分别如下，可能是二叉搜索树的是（\quad）。
\begin{enumerate}
  \item $5,3,1,2,4$
  \item $1,2,3,4,5$
  \item $5,3,4,2,1$
  \item $1,3,5,4,2$
\end{enumerate}
\begin{solution}
选 B。二叉搜索树的中序遍历序列应为非递减序列。
\end{solution}
\end{question}

\begin{question}
线性表有 $2n\;(n>100)$ 个元素，以下操作中，哪一项在单链表上实现比在顺序表上实现效率更高？（\quad）
\begin{enumerate}
  \item 在表中最后一个元素的后面插入一个新元素
  \item 顺序输出表中的前 $i$ 个元素
  \item 交换表中第 $i$ 个元素和第 $n+i$ 个元素的值（$i=0,\ldots,n-1$）
  \item 删除表中第 $1$ 个元素
\end{enumerate}
\begin{solution}
选 D。顺序表删除首元素需移动大量后继元素；单链表删除首结点只需修改头指针。
\end{solution}
\end{question}

\begin{question}
用数组 $O[n]$ 实现循环队列。设队头的前一个元素下标为 $f$，队尾下标为 $r$，则队列中元素总数为（\quad）。
\begin{enumerate}
  \item $r-f$
  \item $(n+f-r)\bmod n$
  \item $n+r-f$
  \item $(n+r-f)\bmod n$
\end{enumerate}
\begin{solution}
选 D。循环队列中从 $f$ 的后继走到 $r$ 的元素个数为 $(r-f+n)\bmod n$。
\end{solution}
\end{question}

\begin{question}
一个拥有 $21$ 条边的非连通无向图，至少包含（\quad）个顶点。
\begin{enumerate}
  \item $5$
  \item $6$
  \item $7$
  \item $8$
\end{enumerate}
\begin{solution}
选 D。$v$ 个顶点的非连通无向简单图边数至多为 $\binom{v-1}{2}$。当 $v=7$ 时最多 $15$ 条边，不足 $21$；当 $v=8$ 时最多 $21$ 条边。
\end{solution}
\end{question}

\begin{question}
若使用分治算法解决某问题，每次把规模为 $n$ 的问题分解为 $2$ 个规模为 $n/2$ 的子问题，并用 $O(n)$ 时间合并两个子问题结果，则总时间复杂度为（\quad）。
\begin{enumerate}
  \item $O(n)$
  \item $O(n\log n)$
  \item $O(n^2)$
  \item $O(n^2\log n)$
\end{enumerate}
\begin{solution}
选 B。递推式为 $T(n)=2T(n/2)+O(n)$，由主定理得 $T(n)=O(n\log n)$。
\end{solution}
\end{question}

\begin{question}
有 $n$ 个顶点、$m$ 条边的无向图，下列说法中错误的是（\quad）。
\begin{enumerate}
  \item 采用邻接表存储，空间复杂度为 $O(n+m)$
  \item 若为稠密图，更倾向于采用邻接矩阵存储
  \item 采用邻接表存储，删除一个顶点的最坏情况时间复杂度为 $O(n+m)$
  \item 若为稀疏图，深度优先周游的时间复杂度低于广度优先周游的时间复杂度
\end{enumerate}
\begin{solution}
选 D。邻接表存储下 DFS 与 BFS 的时间复杂度均为 $O(n+m)$。
\end{solution}
\end{question}

\begin{question}
利用栈将表达式 $3*2\string^(4+2*2-6*3)-5$（其中 \texttt{\string^} 为乘幂）转换为后缀表达式。扫描到 $6$ 时，运算符栈为（\quad）。
\begin{enumerate}
  \item \texttt{*\string^(+-}
  \item \texttt{*\string^}
  \item \texttt{*\string^(+}
  \item \texttt{*\string^(-}
\end{enumerate}
\begin{solution}
选 D。扫描到 $6$ 前，括号内的 $+$ 与 $*$ 已因遇到 $-$ 出栈，栈中自底向上为 \texttt{*\string^(-}。
\end{solution}
\end{question}

\begin{question}
定义一棵没有 $1$ 度结点的二叉树为满二叉树。对于一棵包含 $k$ 个结点的满二叉树，其中叶子结点的个数为（\quad）。
\begin{enumerate}
  \item $\lfloor k/2\rfloor$
  \item $\lfloor k/2\rfloor-1$
  \item $\lfloor k/2\rfloor+1$
  \item 以上三个都有可能
\end{enumerate}
\begin{solution}
选 C。满二叉树中若分支结点数为 $m$，叶子数为 $m+1$，总数 $k=2m+1$，故叶子数为 $(k+1)/2=\lfloor k/2\rfloor+1$。
\end{solution}
\end{question}

\textbf{二、判断题（10 分，每小题 1 分；对填 Y，错填 N）}
\begin{question}
判断下列命题正误。
\begin{enumerate}
  \item 分治法和动态规划法都适用于将问题分解为规模较小的子问题的思想。
  \item 在链表中查找和删除一个元素的时间复杂度都是 $O(1)$。
  \item 相比于顺序存储，完全二叉树更适合用链式表示存储。
  \item 通常不能仅通过一棵二叉树的前序遍历结点序列和后序遍历结点序列确定这棵二叉树的完整结构。
  \item 二叉搜索树一定是满二叉树。
  \item 归并排序算法一定比简单插入排序算法的执行效率高。
  \item 在待排序元素为正序情况下，直接插入排序可能比快速排序的时间复杂度更小。
  \item 一个有 $n$ 个结点的无向图，最少有一个连通分量。
  \item 图的遍历算法（如深度优先搜索 DFS 和广度优先搜索 BFS）都只能用于无向图。
  \item 若连通无向图的边权值互不相同，其最小生成树只有一个。
\end{enumerate}
\begin{solution}
依次为：Y，N，N，Y，N，N，Y，Y，N，Y。
\end{solution}
\end{question}

\textbf{三、填空题（20 分，每题 2 分）}
\begin{question}
根据题意填空。
\begin{enumerate}
  \item 若经常需要在线性表头部进行插入和删除操作，最合适的存储结构是链式存储；若需要随机存取，最合适的存储结构是顺序存储。
  \item 顺序表中 $a,b,c,d,e,f$ 的查找概率分别为 $0.12,0.25,0.23,0.20,0.05,0.15$。为了使查找成功的平均比较次数最少，数据元素的存放顺序应为 $b,c,d,f,a,e$。
  \item 含 $120$ 个结点的完全二叉树共有 $7$ 层，有 $60$ 个叶子结点。
  \item 已知一棵树的中序遍历序列为 \texttt{DBGEACF}、后序遍历序列为 \texttt{DGEBFCA}，其前序遍历结果为 \texttt{ABDEGCF}。
  \item 字符 $a,b,c,d$ 的出现次数分别为 $1000,2000,6000,1000$，采用 Huffman 编码时，该电文总长度为 $16000$ 比特。
  \item 对关键字序列 $45,80,55,40,42,85$ 从最后一个非叶子结点开始调整建最大堆，得到的初始最大堆为 $85,80,55,40,42,45$。
  \item 一棵 $m$ 叉树中度数为 $i$ 的结点数为 $N_i$，则终端结点数为 $1+\sum_{i=2}^{m}(i-1)N_i$。
  \item 散列表长 $m=15$，散列函数 $H(\mathrm{key})=\mathrm{key}\%11$，已有地址 $\mathrm{addr}(16)=5$、$\mathrm{addr}(37)=4$、$\mathrm{addr}(50)=6$、$\mathrm{addr}(70)=7$。若用线性探查处理冲突，关键字 $38$ 的地址为 $8$。
  \item 森林 $F$ 有 $4$ 棵树，结点数分别为 $10,9,11,7$。转成一棵二叉树后，其根结点右子树中有 $27$ 个结点。
  \item 对无向图，若边数 $m\ge n$（$n$ 为结点数），则该图一定存在回路。
\end{enumerate}
\begin{solution}
见各小题填空。
\end{solution}
\end{question}

\textbf{四、简答题（3 题，共 14 分）}
\begin{question}
对序列 $(H,E,B,L,G,A,F,J,I,C,D,K)$ 按字母升序排序，求：
\begin{enumerate}
  \item 冒泡排序第一趟交换结束后的结果；
  \item 二路归并排序第一趟归并后的结果；
  \item 初始步长为 $4$ 的 Shell 排序第一趟后的结果。
\end{enumerate}
\begin{solution}
\begin{enumerate}
  \item 冒泡排序第一趟：$(E,B,H,G,A,F,J,I,C,D,K,L)$。
  \item 二路归并第一趟：$(E,H,B,L,A,G,F,J,C,I,D,K)$。
  \item Shell 排序步长为 $4$ 的第一趟：$(G,A,B,J,H,C,D,K,I,E,F,L)$。
\end{enumerate}
\end{solution}
\end{question}

\begin{question}
在题图给出的二叉排序树中删除结点 $37$。要求先寻找该结点左子树上的最大者 $r$，并利用 $r$ 辅助删除该结点。请画出删除后的二叉排序树结构。
\pyfig[0.38\linewidth]{figures/数据结构与算法B-24春-期末-fig01.png}{}

\begin{solution}
左子树最大结点为 $32$。删除后以 $32$ 替代原根 $37$；左子树根为 $24$，其左孩子为 $7$、右孩子为 $30$，$7$ 的左孩子为 $2$；右子树保持为根 $42$，其左孩子为 $40$、右孩子为 $42$，该右孩子的右孩子为 $120$。
\end{solution}
\end{question}

\begin{question}
Dr. Stranger 的电脑感染了一种病毒。该病毒会把文档中的每种字母固定替换为另一种字母，且互不相同，可以看作字母集合上的双射函数。已知字母集合 $X=\{a,b,c,d,e\}$，文档感染后内容为
\[
D'=\{\texttt{cedbdac},\texttt{cac},\texttt{ecd},\texttt{dca},\texttt{aba},\texttt{bac}\}.
\]
要求破解病毒规则 $f$ 并还原原文档 $D$。
\begin{solution}
按扫描题面仅给出密文集合和“每个字母替换为另一不同字母”的条件，替换规则并不能唯一确定：任意一个 $X$ 上的无不动点置换都能对应某个原文集合。若考试另有词典、明文样例或上下文约束，应再据此确定 $f$；在当前题面下，严格结论是条件不足，原文不唯一。
\end{solution}
\end{question}

\textbf{五、算法填空（4 题，共 26 分）}
\begin{question}[只带尾指针的循环单链表]
程序描述只带尾指针的循环单链表操作：先把 $n$ 个整数依次插入链表头部，然后删除链表尾部 $m$ 个元素，再在链表前部插入 $k$ 个整数。补全程序空白。
\begin{solution}
空白处依次为：
\begin{center}
\small
\begin{tabular}{@{}r l@{}}
(1)&\texttt{nd->next = tail->next},\\
(2)&\texttt{tail->next == tail},\\
(3)&\texttt{ptr->next != tail},\\
(4)&\texttt{tail = ptr},\\
(5)&\texttt{ptr != tail->next}.
\end{tabular}
\end{center}
\end{solution}
\end{question}

\begin{question}[二叉树宽度]
给定一棵二叉树，宽度定义为结点最多的那一层的结点数。输入每个结点的左右儿子编号，若为 $-1$ 表示没有相应儿子，补全程序求二叉树宽度。
\begin{solution}
空白处依次为：
\begin{center}
\small
\begin{tabular}{@{}r l@{}}
(1)&\texttt{traversal(root->left, level + 1)},\\
(2)&\texttt{traversal(root->right, level + 1)},\\
(3)&\texttt{nodes[L].parent = nd},\\
(4)&\texttt{nodes[R].parent = nd},\\
(5)&\texttt{nodes[i].parent == -1}.
\end{tabular}
\end{center}
这里用 \texttt{parent} 临时记录父结点编号，未被记录为任何结点儿子的结点即为根。
\end{solution}
\end{question}

\begin{question}[Prim 算法]
用 Prim 算法求无向图最小生成树。图用邻接矩阵 $G$ 存储，$G[u][v]=\texttt{INF}$ 表示无边，顶点编号从 $0$ 开始。补全程序，使其返回最小生成树的边列表。
\begin{solution}
空白处依次为：
\begin{center}
\small
\begin{tabular}{@{}r l@{}}
(1)&\texttt{doneNum < n},\\
(2)&\texttt{dist[i] < minDist},\\
(3)&\texttt{used[x] = true},\\
(4)&\texttt{make\_pair(x, prev[x])},\\
(5)&\texttt{G[x][v] != INF \&\& G[x][v] < dist[v]},\\
(6)&\texttt{prev[v] = x}.
\end{tabular}
\end{center}
\end{solution}
\end{question}

\begin{question}[连通分量大小]
补全程序，计算一个无向图中所有连通分量的结点数。输入为图的结点数、边数和边列表，输出每个连通分量的结点数。
\begin{solution}
空白处依次为：
\begin{center}
\small
\begin{tabular}{@{}r l@{}}
(1)&\texttt{!visited[i]},\\
(2)&\texttt{dfs(i, visited, graph, n)},\\
(3)&\texttt{graph[u].push\_back(v)},\\
(4)&\texttt{graph[v].push\_back(u)},\\
(5)&\texttt{!visited[i]},\\
(6)&\texttt{dfs(i, visited, graph, n)}.
\end{tabular}
\end{center}
\end{solution}
\end{question}
